\chapter{生物性有害因素所致职业性损害}
\setcounter{chapter}{7}
\chapterintro{
\textbf{【教学大纲要求】}

本章为\imp{自学}内容。包括职业性炭疽、布鲁氏菌病、职业性森林脑炎（蜱传脑炎）等生物性有害因素所致职业性损害。
}

\section{职业性生物性有害因素概述}

\begin{defbox}
\textbf{生物性有害因素}：生产原料和生产环境中存在的对职业人群健康有害的致病微生物、寄生虫、昆虫和其他动植物及其所产生的生物活性物质。常见例子：动物皮毛上的炭疽芽胞杆菌、布鲁氏菌，蜱媒森林脑炎病毒、支原体、衣原体、钩端螺旋体等。
\end{defbox}

\section{职业性炭疽}

\begin{keybox}
\textbf{炭疽（Anthrax）}：由炭疽芽胞杆菌引起的人畜共患传染病，主要发生于牛、羊、马等食草动物；人主要通过接触患病动物/制品或吸入炭疽芽胞而感染。法定管理：乙类传染病，其中\keyword{肺炭疽按甲类传染病管理}。职业性炭疽是\keyword{国家法定职业病}。

\textbf{病原学与致病机制：}
\begin{itemize}[leftmargin=*]
    \item 主要致病物质：荚膜和外毒素
    \item 荚膜：由D-谷氨酸多肽组成，具抗吞噬作用
    \item 外毒素：由保护性抗原（PA）、致死因子（LF）、水肿因子（EF）三种蛋白组成
    \item 皮肤感染：炭疽杆菌侵入皮下组织→释放毒素→细胞水肿和组织坏死→原发性皮肤炭疽
    \item 呼吸道感染：吸入炭疽芽孢→出血性肺炎和肺门淋巴结炎
    \item 消化道感染：食入→急性肠炎（出血性炎症、组织水肿坏死）
    \item 全身中毒：损害微血管内皮细胞→DIC、感染性休克→多器官功能衰竭
\end{itemize}

\textbf{流行病学：}传染源主要为染疫食草动物；人群普遍易感，农牧民、饲养员、屠宰/皮毛加工人员、兽医、检疫人员等职业人群感染风险更高；感染后可获得持久免疫力。

\textbf{防治原则：}肺炭疽按甲类传染病管理，发现后2小时内网络直报；患者就地隔离治疗；皮肤炭疽按乙类传染病隔离管理，皮损原则上不做病灶切除和引流；肺炭疽执行标准预防+空气隔离+飞沫隔离。
\end{keybox}

\section{布鲁氏菌病}

布鲁氏菌有荚膜，可产生透明质酸酶和过氧化氢酶，侵袭力强，能通过完整皮肤和黏膜进入宿主体内。重要致病因子：IV型分泌系统（T4SS）和脂多糖（LPS）。胞内寄生特点：在局部淋巴结生长繁殖→突破屏障侵入血液→菌血症→细菌反复入血→波浪热。

\section{职业性森林脑炎（蜱传脑炎）}

森林脑炎病毒（森脑病毒）仅存在于自然疫源地，寄生于啮齿类动物及鸟类血液中，通过蜱类媒介传播。蜱类既是传播媒介，也是病毒的长期宿主。人类主要通过蜱虫叮咬或消化途径感染，人作为宿主的传染性极低。职业性森林脑炎是劳动者在森林地区从事职业活动时因被蜱叮咬感染引发的中枢神经系统急性病毒性传染病，属于国家法定职业病。

\newpage

% =====================================================================
%  CHAPTER 8: 其他职业病
% =====================================================================
\newpage